Extending baseline training lowered speed distribution error from 709.73 to 51.64 m/s and the fraction of generated points outside the training bounds from 88.4\% to 8.2\%. Short trips remained particularly difficult. Across four equally sized groups ordered by duration, speed $W_1$ ranged from 139.55 m/s in the shortest group to 8.94 m/s in the longest. In the shortest group, successive generated points had a median separation of 31.35 m, compared with 4.52 m for observed points.

The table compares two training seeds at 8,000 updates. In each seed, speed supervision alone reduced speed and length error while increasing spatial density divergence. The absolute gain depended strongly on the baseline fit. Class conditioning alone increased speed error in both seeds. The combined model yielded larger motion errors than either baseline. A time-varying target thus improved particular motion statistics under these settings. The class embedding supplied no repeatable improvement on those endpoints. An independent vehicle population is needed to estimate the size of the effect beyond these development requests.
